\documentclass[11pt]{article}

\usepackage[a4paper,margin=2.2cm]{geometry}
\usepackage[utf8]{inputenc}
\usepackage[T1]{fontenc}
\usepackage{lmodern}
\usepackage{amsmath, amssymb, amsthm, mathtools}
\theoremstyle{definition}
\newtheorem*{definition}{Definition}
\newtheorem*{remark}{Remark}
\usepackage{microtype}
\usepackage{enumitem}
\usepackage{xcolor}
\usepackage{hyperref}
\usepackage{titlesec}
\usepackage{fancyhdr}
\usepackage{booktabs}

\definecolor{accent}{HTML}{1F3A68}
\definecolor{warn}{HTML}{8B2A2A}
\definecolor{soft}{HTML}{E8EEF5}

\hypersetup{colorlinks=true, linkcolor=accent, urlcolor=accent}

\titleformat{\section}{\large\bfseries\color{accent}}{\thesection.}{0.6em}{}
\titleformat{\subsection}{\normalsize\bfseries\color{accent}}{}{0em}{}

\pagestyle{fancy}
\fancyhf{}
\fancyhead[L]{\small\textit{Speaker script --- Bootstrap Inference in the Presence of Bias}}
\fancyhead[R]{\small\thepage}
\renewcommand{\headrulewidth}{0.3pt}

% ===== Macros =====
\newcommand{\R}{\mathbb{R}}
\newcommand{\E}{\mathbb{E}}
\newcommand{\PR}{\mathbb{P}}
\newcommand{\Pstar}{P^{*}}
\newcommand{\Pstst}{P^{**}}
\newcommand{\dto}{\xrightarrow{d}}
\newcommand{\pto}{\xrightarrow{p}}
\newcommand{\dstarto}{\xrightarrow{d^{*}}}
\newcommand{\dstarstarto}{\xrightarrow{d^{**}}}
\newcommand{\Tn}{T_{n}}
\newcommand{\Tns}{T_{n}^{*}}
\newcommand{\Tnss}{T_{n}^{**}}
\newcommand{\Bn}{B_{n}}
\newcommand{\Bnh}{\hat{B}_{n}}
\newcommand{\Bnhs}{\hat{B}_{n}^{*}}
\newcommand{\thn}{\hat{\theta}_{n}}
\newcommand{\thns}{\hat{\theta}_{n}^{*}}
\newcommand{\pnh}{\hat{p}_{n}}
\newcommand{\pnhs}{\hat{p}_{n}^{*}}
\newcommand{\pnt}{\tilde{p}_{n}}
\newcommand{\Hnh}{\hat{H}_{n}}
\newcommand{\Lnh}{\hat{L}_{n}}
\newcommand{\Dn}{D_{n}}
\newcommand{\Dns}{D_{n}^{*}}
\newcommand{\Dnss}{D_{n}^{**}}

\newcommand{\slide}[2]{%
  \medskip
  \noindent\colorbox{soft}{\parbox{\linewidth}{\textbf{Slide #1.} \textit{#2}}}
  \medskip
}
\newcommand{\say}[1]{\textbf{Say:} #1}
\newcommand{\why}[1]{\textbf{Why this slide:} #1}
\newcommand{\trans}[1]{\textbf{Transition:} #1}
\newcommand{\Q}[2]{\noindent\textcolor{warn}{\textbf{Q:}} #1\\ \textcolor{accent}{\textbf{A:}} #2\par\smallskip}

\title{\color{accent}\bfseries Speaker Script\\
\large Bootstrap Inference in the Presence of Bias\\
\normalsize\mdseries Cavaliere, Gonçalves, Nielsen \& Zanelli (JASA, 2024)}
\author{Jordi Torres --- TSE Econometrics Reading Group}
\date{}

\begin{document}
\maketitle
\thispagestyle{empty}

\noindent\textit{This document is a slide-by-slide companion to
\texttt{presentation\_bootstrapbias.tex}. For each slide it contains:
(i) a tight script, (ii) the notation introduced, (iii) the transition into
the next slide, and (iv) anticipated aggressive questions with prepared answers.
A target of $\sim$45 minutes implies roughly 1--1.5 minutes per slide.}

\bigskip
\hrule
\bigskip

% ====================================================================
\section{Master cheat-sheet of notation}

Define each symbol \emph{once} below and use it throughout. After the relevant
slide introduces it, do \emph{not} re-define.

\subsection*{Population / data}
\begin{description}[leftmargin=2.6cm,style=nextline]
  \item[$\theta\in\R$] scalar parameter of interest.
  \item[$\Dn$] original sample of size $n$.
  \item[$\Dns$] bootstrap sample drawn from a bootstrap DGP \emph{conditional} on $\Dn$.
\end{description}

\subsection*{Statistics}
\begin{description}[leftmargin=2.6cm,style=nextline]
  \item[$\thn$] estimator of $\theta$ based on $\Dn$.
  \item[$\thns$] bootstrap version of $\thn$, computed from $\Dns$.
  \item[$g(n)$] convergence rate ($n^{1/2}$, $(nh)^{1/2}$, etc.).
  \item[$\Tn$] $:= g(n)(\thn-\theta)$, the centred and rescaled statistic.
  \item[$\Tns$] $:= g(n)(\thns-\thn)$, the bootstrap version
    (note centring at $\thn$, not $\theta$).
\end{description}

\subsection*{Bias terms (the heart of the paper)}
\begin{description}[leftmargin=2.6cm,style=nextline]
  \item[$\Bn$] deterministic asymptotic bias sequence: $\Bn\to B$.
  \item[$B$] limiting bias of $\Tn$ (constant; \emph{may not be consistently estimable}).
  \item[$\Bnh$] \emph{implicit} bootstrap bias --- $\Dn$-measurable, defined so that
    $\Tns-\Bnh$ has the same asymptotic noise as $\Tn-\Bn$. Crucially \emph{not} an estimator the researcher chooses; it is whatever bias the bootstrap DGP produces.
  \item[$\Bnhs$] second-level analogue used in the double bootstrap.
\end{description}

\subsection*{Limit objects}
\begin{description}[leftmargin=2.6cm,style=nextline]
  \item[$\xi_1$] limit of $\Tn-\Bn$, centred at zero, cdf $G$ continuous.
  \item[$\xi_2$] limit of $\Bnh-\Bn$, centred at zero. \emph{Nonzero $\xi_2$ is the failure.}
  \item[$G$] cdf of $\xi_1$.
  \item[$F$] cdf of $\xi_1-\xi_2$ (continuous by assumption).
  \item[$H$] $:= F\circ G^{-1}$, asymptotic cdf of $\pnh$. Continuous, free of $B$.
\end{description}

\subsection*{Bootstrap p-values}
\begin{description}[leftmargin=2.6cm,style=nextline]
  \item[$\pnh$] $:= \Pstar(\Tns\le \Tn)$, standard bootstrap $p$-value.
  \item[$\pnhs$] second-level $p$-value $\Pstst(\Tnss\le\Tns)$ for the double bootstrap.
  \item[$\Hnh$] estimator of $H$ (plug-in or double bootstrap).
  \item[$\pnt$] $:= \Hnh(\pnh)$, the prepivoted (modified) $p$-value.
\end{description}

\subsection*{Modes of convergence}
\begin{description}[leftmargin=2.6cm,style=nextline]
  \item[$\pto$] convergence in probability under $P$.
  \item[$\dto$] convergence in distribution under $P$.
  \item[$\Pstar$] $\!:=P(\cdot\mid\Dn)$, bootstrap probability measure (random in $\Dn$).
  \item[$\pto_p$ \,(written $\xrightarrow{p^{*}}_{p}$)] $Y_n^*\pto_p 0$ means
    $\Pstar(|Y_n^*|>\varepsilon)\pto 0$ for every $\varepsilon>0$. (Convergence to $0$
    in bootstrap probability \emph{in probability under $P$}.)
  \item[$\dstarto_p$] $Y_n^*\dstarto_p\xi$ means
    $\Pstar(Y_n^*\le u)\pto P(\xi\le u)$ at every continuity point.
\end{description}

\textbf{Read aloud only when first introduced.} After the bootstrap-statistic slide,
just say ``in bootstrap probability'' or ``conditionally on the data.''

\bigskip
\hrule
\bigskip

% ====================================================================
\section{Part 1 --- Introduction (slides 1--8)}

\slide{1}{Title slide.}
\say{Today I am presenting Cavaliere, Gonçalves, Nielsen \& Zanelli, JASA 2024.
The paper is about how to do valid bootstrap inference when the estimator has an
asymptotic bias the bootstrap cannot replicate. Before going into details, let
me place the paper in our reading group.}

\trans{Move to the roadmap.}

% ----------------
\slide{2}{Roadmap.}
\say{Five blocks: a short \emph{introduction} that fixes notation and states the
problem; three \emph{motivating examples}; the \emph{general theory} ---
Theorem~3.1 and prepivoting --- which is the core of the talk; the
\emph{examples revisited} where I show how the fix is implemented; and a
short \emph{wrap-up}.}

\trans{Begin with the framing slide.}

% ----------------
\slide{3}{Placing the paper in the reading group.}
\say{In the reading group so far we have seen two roles for the bootstrap.
First, computing variances and full sampling distributions when analytical
formulas are unavailable. Second, asymptotic refinements --- Beran 1987 and
Hall 1992 --- where the bootstrap can be more accurate than the first-order
normal approximation.

Today's paper is different. It asks: what if the estimator has an asymptotic
bias $B$ that the bootstrap simply cannot replicate? Then the standard
bootstrap is invalid, even asymptotically. Cavaliere et al.\ show that valid
inference can be restored \emph{without estimating the bias} by a transformation
of the bootstrap $p$-value, the prepivoting idea of Beran 1987, 1988.}

\why{This is the only narrative slide in the deck. From here on we go technical.}

\trans{Now I make the setup precise.}

\textbf{Anticipated questions:}
\Q{Why is bias different from variance? Cant we just bootstrap the bias?}
{Variance is a property of the sampling distribution that the bootstrap
mimics by construction. Bias is the gap between the population parameter and
the limit of the estimator; the bootstrap replicates the law of $\Tns$ \emph{around
$\thn$}, not the law of $\Tn$ around $\theta$, so bias does not transfer.}
\Q{Is this related to the failure of the bootstrap under the null in
unit-root or moment-condition problems?}
{Conceptually yes --- in all those cases the bootstrap measure is random
in the limit. The unifying language here is condition (1.3) ahead.}

% ----------------
\slide{4}{Setup --- the original statistic (1.1).}
\say{Let $\theta\in\R$ be a scalar parameter, $\thn$ an estimator with rate
$g(n)\to\infty$. Define $\Tn:=g(n)(\thn-\theta)$. Equation (1.1) says
$\Tn\dto B+\xi_1$, with $\xi_1$ centred at zero and continuous cdf $G$.

The point: $B$ is an \emph{asymptotic bias}. In every example we will see, $B$
is \emph{not} consistently estimable. The typical setting is $g(n)=n^{1/2}$ and
$\xi_1$ Gaussian, but the paper allows non-Gaussian $\xi_1$.

Why does classical inference fail? Because using quantiles of $\xi_1$ to build
confidence intervals or tests \emph{ignores} $B$, so coverage is wrong.

Why not just bias-correct analytically? Either $B$ cannot be estimated (model
averaging, ridge with local-to-zero parameter), or it can be estimated but only
poorly (nonparametric regression: $\beta''(x)$ requires a second bandwidth).
And as the box notes, the bootstrap also \emph{cannot} solve the bias problem on
its own when $B$ is not estimable.}

\textbf{Notation introduced:} $\theta,\thn,\Tn,\Bn,B,\xi_1,G,g(n)$.

\trans{Now, what does the bootstrap give us?}

\textbf{Anticipated questions:}
\Q{What is ``not consistently estimable'' precisely? Lots of biases are slow but estimable.}
{Take model averaging: bias is $Q\cdot c$ where $\delta=cn^{-1/2}$. The
parameter $c$ is unidentified --- $\hat\delta_n=O_p(n^{-1/2})$ converges in
\emph{distribution}, not in probability, when scaled by $n^{1/2}$. So no estimator
can be consistent for $c$. Same logic for ridge.}
\Q{Doesnt Hall (1992) handle this with Edgeworth refinements?}
{No --- those refinements assume the bootstrap is first-order correct and then
chase higher-order terms. Here it fails at first order.}

% ----------------
\slide{5}{Setup --- the bootstrap statistic (1.2).}
\say{$\Tns:=g(n)(\thns-\thn)$ --- centred at $\thn$ because, in the bootstrap
world, $\thn$ plays the role of the ``true'' parameter (the only feasible
choice; $\theta$ is unknown). Equation (1.2): $\Tns-\Bnh\dstarto_p\xi_1$
--- the bootstrap reproduces the noise $\xi_1$ correctly.

Critically, $\Bnh$ is \emph{implicit}: it is whatever asymptotic bias the
bootstrap DGP \emph{produces}, not something the researcher chooses or estimates.
You can think of $\Bnh$ as ``the bootstrap's own implicit estimator of $B$ ---
even though no one wrote it down.''

If $\Bnh-B=o_p(1)$ --- i.e.\ the implicit estimator is consistent --- the
bootstrap is valid in the usual sense. The whole question of the paper is:
what happens when no such consistent implicit estimator can exist?

The notation block on the right defines two modes of bootstrap convergence I
will use repeatedly.}

\textbf{Notation introduced:} $\Tns,\Bnh,\Dn,\Dns,\Pstar,\dstarto_p,\pto_p$.

\trans{Now I state the failure mode.}

\textbf{Anticipated questions:}
\Q{What does ``implicit'' mean? Where is $\Bnh$ defined operationally?}
{Operationally $\Bnh$ is whatever sequence makes (1.2) hold --- it is a
property of the bootstrap DGP, not a choice. In the model averaging FRB
example we will see $\Bnh = Q_n\, n^{1/2}\hat\delta_n$.}
\Q{Why centre $\Tns$ at $\thn$ rather than $\theta$?}
{Because $\theta$ is unknown; centring at $\thn$ is the only feasible choice.
This is also why the bootstrap distribution is conditional on $\Dn$.}

% ----------------
\slide{6}{The failure mode (1.3).}
\say{Condition (1.3): $\Bnh-B\dto\xi_2\not\equiv 0$, jointly with (1.1).
This is the failure mode --- and the key idea worth dwelling on.

\emph{Why does this happen?} $\Bnh$ is built from the data, and in the
problems we care about, the data noise sits at the \emph{same rate} as
the bias $B$. So no estimator of $B$ from this data can be consistent ---
it inherits an irreducible $O_p(1)$ noise of the same order as the thing
it is trying to estimate. The bootstrap's implicit $\Bnh$ is no exception.

\emph{Two noises, not one.} $\xi_1$ is the sampling noise in $\Tn$ ---
the bootstrap reproduces this. $\xi_2$ is the bootstrap's failure to
pin down the bias --- this is an \emph{extra} source of randomness that
the bootstrap cannot subtract out, because it \emph{is} the bootstrap's
own error.

\emph{Consequences.} The bootstrap distribution of $\Tns$ conditional on
$\Dn$ is itself \emph{random in the limit} --- it sits at a random shift
$\Bnh$, and that randomness $\xi_2$ never washes out. So $\Tns$ and $\Tn$
have \emph{different} asymptotic laws (one centred at $\Bnh$, the other at
$\Bn$, differing by $\xi_2$). Hence $\pnh:=\Pstar(\Tns\le\Tn)$ is \emph{not}
asymptotically uniform, and a nominal 5\% test can reject at, say, 15\%.

\emph{Counterfactual to keep in mind:} if $\Bnh-\Bn=o_p(1)$ (so $\xi_2=0$),
then $\Tns$ and $\Tn$ would have the same asymptotic law and $\pnh$ would
be uniform --- standard story. (1.3) is exactly the obstruction.

The goal of the paper: restore valid inference without trying to estimate
$B$ --- by working with the bootstrap $p$-value $\pnh$ directly.}

\textbf{Notation introduced:} $\xi_2,\,\pnh$.

\trans{The fix is on the next slide.}

\textbf{Anticipated questions:}
\Q{Random in the limit --- is this related to Cavaliere--Georgiev (2020) or to
random measures?}
{Yes, exactly. CG 2020 showed randomness of the bootstrap measure does not by
itself prevent validity. Here it does, because the asymptotic bias is also random.}
\Q{Could we just use a bigger bootstrap, $m$-out-of-$n$, to fix this?}
{Often $m$-out-of-$n$ \emph{kills} the bias in the bootstrap world, giving
$\Bnh\pto 0$ and $\xi_2=-B$ --- a constant. By Remark 3.1 ahead, that case
is \emph{worse}, not better: the distortion now depends on $B$, so prepivoting
no longer works either.}

% ----------------
\slide{7}{A way out --- prepivoting.}
\say{Beran's idea, 1987--88. Even though $\pnh$ is not uniform, its
asymptotic cdf $H$ is going to be \emph{continuous} and, more importantly,
\emph{free of the unknown bias $B$}. If so, the probability integral transform
gives $H(\pnh)\dto U[0,1]$. Replace $\pnh$ by $\pnt:=\Hnh(\pnh)$, with $\Hnh$
uniformly consistent for $H$.

Two implementations: \emph{plug-in} when $H$ depends on a finite-dimensional
parameter, and \emph{double bootstrap} when it does not. Beran proposed this
for refinements; here it rescues a broken bootstrap.}

\textbf{Notation introduced:} $H,\Hnh,\pnt$.

\trans{Three motivating examples.}

\textbf{Anticipated questions:}
\Q{Why should $H$ be free of $B$ in general?}
{This is the content of Theorem~3.1, which I will prove. The reason is that
both numerator and denominator in $H=F\circ G^{-1}$ are cdfs of \emph{centred}
limits $\xi_1$ and $\xi_1-\xi_2$, and $B$ is the centring that has been
removed by definition.}
\Q{Is prepivoting a refinement or a fix?}
{Beran used it for refinements --- improving second-order accuracy. Here it
fixes a first-order failure. Mechanically the operation is the same.}

% ====================================================================
\section{Part 2 --- Motivating examples (slides 9--14)}

\slide{9}{Three motivating examples (lead-in).}
\say{For each example I will check four things in order: the model and
parameter; whether (1.1) holds; whether (1.2) holds; and whether (1.3) holds.
The third one --- (1.3) --- is the key to understanding why the bootstrap
fails.}

\trans{Example 1, model averaging.}

% ----------------
\slide{10}{Example 1 --- Ridge regression, setup.}
\say{Linear regression with $p$ regressors. Test a linear restriction
$g'\theta=r$. Ridge estimator with shrinkage $c_n$ added to the diagonal
of $S_{xx}$. The local-to-zero device is $\theta=\delta n^{-1/2}$ and
$c_n/n\to c_0>0$. As before, this is the regime where shrinkage bias
matters.}

\textbf{Anticipated questions:}
\Q{$c_n\sim n$ is a strong tuning. Is that standard?}
{It is the regime studied by Fu and Knight (2000) and Hjort and
Claeskens (2003). The point is to keep shrinkage at the same order as noise.
Other rates lead to either no shrinkage in the limit or to $\theta=0$ being
imposed asymptotically.}

% ----------------
\slide{11}{Example 1 --- Why (1.1)--(1.3) hold.}
\say{(1.1) --- bias survives in the limit. From
$\tilde S_{xx}^{-1}S_{xx}-I=-(c_n/n)\tilde S_{xx}^{-1}$,
$\Tn$ splits into a shrinkage piece (limit $B \propto c_0\,\delta$) and a
Gaussian noise piece (limit $\xi_1\sim N(0,v^2)$). The bias depends on the
local-to-zero $\delta$, which is \emph{not consistently estimable} because
$n^{1/2}\hat\theta_n=\delta+O_p(1)$ converges in \emph{distribution}, not in
probability.

(1.2) --- bootstrap captures the noise. Pairs bootstrap, \emph{centred at OLS},
not at ridge. Reason: the bootstrap world's ``true parameter'' must satisfy
the score condition $E^*[x_t^*\varepsilon_t^*]=0$; OLS does, ridge does not.
Same recentring trick as in the GMM bootstrap (Horowitz, Hall--Horowitz).

(1.3) --- and here is the key point. The implicit bootstrap bias is
$\Bnh = -(c_n/n)g'\tilde S_{xx}^{-1}\,n^{1/2}\hat\theta_n$. Subtracting
$\Bn$:
\[
  \Bnh-\Bn = -\tfrac{c_n}{n}g'\tilde S_{xx}^{-1}\,n^{1/2}(\hat\theta_n-\theta).
\]
$c_n/n\to c_0>0$ and $n^{1/2}(\hat\theta_n-\theta) = O_p(1)$ \emph{not}
$o_p(1)$ --- OLS is only $n^{1/2}$-consistent, so its noise does not vanish at
the relevant scale. Multiplying $O_p(1)$ noise by a non-degenerate matrix
gives $O_p(1)$. Hence $\xi_2\neq 0$.

\emph{Tie to the general story:} the bootstrap is implicitly trying to
estimate $B$ via $\Bnh$, and $\Bnh$ inherits the OLS noise --- which is the
\emph{same order} as the bias itself. That is exactly why the bias is
``random in the limit'' in this example.}

\trans{Example 2, nonparametric.}

\textbf{Anticipated questions:}
\Q{Why pairs and not residual bootstrap?}
{Residual bootstrap requires homoskedasticity and a correctly specified mean
function. Pairs is robust. The recentring at OLS is the technical price.}
\Q{Wouldnt this also fail under wild bootstrap?}
{Yes. The failure is again about $\hat\theta_n-\theta$ being only
$n^{1/2}$-consistent, not $o(n^{-1/2})$.}

% ----------------
\slide{12}{Example 2 --- Nonparametric regression, setup.}
\say{Nonparametric regression with smooth $\beta(\cdot)$, fixed design.
Nadaraya--Watson with MSE-optimal bandwidth $h\sim n^{-1/5}$.

The MSE-optimal $h$ balances squared bias against variance. At this rate,
\emph{the bias does not vanish faster than the noise}; in fact it shows up in
the limit. That is the source of (1.1).}

\textbf{Anticipated questions:}
\Q{Why not undersmooth?}
{Undersmoothing kills bias asymptotically but at the cost of larger variance
and lower power against local alternatives. CCF 2018 document this trade-off.}

% ----------------
\slide{13}{Example 2 --- Why (1.1)--(1.3) hold.}
\say{(1.1) --- bias survives in the limit. Taylor-expand the kernel smoother:
limit bias $B \propto \beta''(x)$, limit noise $\xi_1\sim N(0,v^2)$. The
$\beta''(x)$ is estimable in principle, but only with a second bandwidth ---
unreliable in finite samples (CCF 2018). \emph{Effectively inestimable in
practice.}

(1.2) --- bootstrap captures the noise. Parametric bootstrap with
$y_t^*=\hat\beta_h(x_t)+\varepsilon_t^*$, $\varepsilon_t^*\sim N(0,\hat\sigma_n^2)$.

(1.3) --- here is the key point, with a \emph{different mechanism but the same
moral}. The bootstrap bias replaces $\beta(\cdot)$ with $\hat\beta_h(\cdot)$,
so it depends on $\hat\beta_h''(x)$ instead of $\beta''(x)$. The estimation
error is
\[
  \hat\beta_h''(x)-\beta''(x) = O_p\bigl((nh)^{-1/2}\bigr),
\]
\emph{exactly the same order as the bias itself} at the MSE-optimal bandwidth
(that is what ``MSE-optimal'' means: bias and noise rates are balanced).
Hence $\Bnh-\Bn$ does not vanish, and the bootstrap's implicit estimate of
$B$ inherits irreducible $O_p(1)$ noise.

\emph{Tie to the general story:} same as ridge --- $\Bnh$ is built from data
whose noise sits at the same rate as $B$. The mechanism is different ($O_p(1)$
OLS noise in ridge vs.\ MSE-optimal smoothing here), but the conclusion is
identical: the bootstrap's bias is \emph{random in the limit}.

Cross-example punchline: \textbf{whenever bias and noise are at the same
rate, no estimator of $B$ from the data can be consistent --- the bootstrap
inherits this and (1.3) bites.}}

\trans{Recap, then theory.}

\textbf{Anticipated questions:}
\Q{Couldnt one use a higher-order kernel to remove $B$?}
{Yes, in principle, but at higher variance and at the cost of negative weights
and finite-sample instability. Standard MSE-optimal kernels are second order.}

% ----------------
\slide{14}{Examples --- recap.}
\say{Same diagnostic in both: rate, source of bias, the inestimable
quantity that makes (1.3) bite, the bootstrap scheme. Different mechanisms
($O_p(1)$ nuisance vs.\ MSE-optimal smoothing) but same conclusion. Now we
step back and develop a unified theory.}

\trans{Section 3 of the paper.}

% ====================================================================
\section{Part 3 --- General results (slides 15--24)}

\slide{15}{Recap of notation.}
\say{Quick refresher of the symbols I will use throughout this section. The
bottom block is the most important: $\dstarto_p$ means convergence in
distribution conditional on the data, in probability.

$G$, $F$, $H$ will appear constantly: $G$ is the limit cdf of $\Tn-\Bn$;
$F$ is the limit cdf of $\Tn-\Bnh$, which equals the cdf of $\xi_1-\xi_2$;
$H=F\circ G^{-1}$ will be the asymptotic cdf of $\pnh$.}

\trans{The two assumptions.}

% ----------------
\slide{16}{Assumption 1 --- the original statistic.}
\say{Assumption 1 is just (1.1) restated, plus regularity on $G$:
continuous and strictly increasing on its support.

Continuous, because Polya's theorem upgrades pointwise to uniform convergence
of cdfs --- I will use that in the proof of Theorem 3.1.

Strictly increasing, because we need $G^{-1}$ well-defined to apply the
probability integral transform.

Importantly, $G$ does \emph{not} need to be Gaussian. Stable distributions are
allowed.}

\textbf{Anticipated questions:}
\Q{What if $G$ is discontinuous --- discrete data?}
{Then the framework needs randomisation as in Romano (1990). Out of scope here.}

% ----------------
\slide{17}{Assumption 2 --- the bootstrap side.}
\say{Two parts. (i) is (1.2): the bootstrap captures the noise correctly. (ii)
is the joint convergence of $(\Tn-\Bn,\Bnh-\Bn)$ to $(\xi_1,\xi_2)$, with
$F$ continuous.

Two things to flag. First, $\Bnh$ is again \emph{not} a chosen estimator. It
is whatever the bootstrap DGP produces. Second, $F$ only needs to be
continuous --- not strictly increasing.}

\textbf{Anticipated questions:}
\Q{Joint convergence is strong --- when does it fail?}
{It can fail under non-stationary or near-unit-root settings; see Cavaliere--
Georgiev for examples. In all five examples in the paper, joint convergence
is straightforward because everything is linear in the same score.}

% ----------------
\slide{18}{Benchmark: what if $\Bnh$ \emph{were} consistent?}
\say{Useful warm-up: assume $\Bnh-\Bn=o_p(1)$, so $\xi_2=0$ and $F=G$. Then
$\pnh=\Pstar(\Tns-\Bnh\le\Tn-\Bnh)\to G(\Tn-\Bnh)+o_p(1)\dto G(\xi_1)\sim U[0,1]$.

So when $\xi_2=0$ the bootstrap is uniform: the standard story. The
\emph{non-trivial} case is $\xi_2\neq 0$, where $F\neq G$ and the limit is
$G(F^{-1}(U))\neq U$. That is exactly Theorem~3.1.}

\trans{State and prove Theorem 3.1.}

% ----------------
\slide{19}{Theorem 3.1.}
\say{Under Assumptions 1--2, $\pnh\dto G(F^{-1}(U))$. Three implications.
$\pnh$ is not uniform unless $G=F$. The asymptotic cdf is $H=F\circ G^{-1}$.
And $H$ does not depend on $B$ --- the bias has been subtracted away in both
$F$ and $G$. \emph{This is the algebraic miracle that makes the fix feasible.}}

\textbf{Anticipated questions:}
\Q{Why doesnt $B$ appear in $H$?}
{Both $F$ and $G$ are cdfs of \emph{centred} variables ($\xi_1$ and
$\xi_1-\xi_2$ respectively). Centring removes $B$ by construction. The
information about $B$ is in the location of $\Tn$ itself, not in $H$.}

% ----------------
\slide{20}{Theorem 3.1 --- proof.}
\say{Four steps. \emph{Step 1.} Subtract $\Bnh$ from both sides inside $\Pstar$;
$\Bnh$ is constant under $\Pstar$ since it is $\Dn$-measurable. By Assumption
2(i) and Polya, $\Pstar(\Tns-\Bnh\le u)\to G(u)$ uniformly. So
$\pnh=G(\Tn-\Bnh)+o_p(1)$.

\emph{Step 2.} Decompose $\Tn-\Bnh=(\Tn-\Bn)-(\Bnh-\Bn)$. By Assumption 2(ii)
and CMT, this converges jointly to $\xi_1-\xi_2$.

\emph{Step 3.} $G$ is continuous, so $G(\Tn-\Bnh)\dto G(\xi_1-\xi_2)$.

\emph{Step 4.} $\xi_1-\xi_2$ has continuous cdf $F$, so by the probability
integral transform $\xi_1-\xi_2$ has the same distribution as $F^{-1}(U)$.
Plugging in gives $\pnh\dto G(F^{-1}(U))$. Done.}

\textbf{Anticipated questions:}
\Q{Why is Polya needed?}
{To upgrade pointwise convergence of $\Pstar(\Tns-\Bnh\le u)$ to uniform.
We then evaluate at the random argument $\Tn-\Bnh$ and use Slutsky.
Without continuity of $G$, the substitution does not go through.}
\Q{Where do you actually use joint convergence of (1.1) and (1.3)?}
{In Step 2: the CMT applied to the difference requires joint distributional
convergence, not just marginal.}

% ----------------
\slide{21}{Three remarks.}
\say{Remark 3.1 is the most important. Some bootstrap schemes set $\Bnh\pto 0$ ---
they structurally remove the bias in the bootstrap world. Then $\xi_2=-B$, a
constant. The distortion of $\pnh$ now \emph{depends on $B$}. So
prepivoting cannot fix it: $H$ would depend on the unknown $B$, and you cannot
estimate it. \emph{Lesson: prefer bootstrap schemes that try to capture the bias,
even noisily.}

Remarks 3.2 and 3.3 are corollaries: confidence intervals are similarly
invalid, and the bootstrap conditional distribution is random in the limit.}

\trans{Now use Theorem 3.1 to fix the problem --- Corollary 3.1.}

\textbf{Anticipated questions:}
\Q{Concretely, which schemes give $\Bnh=0$?}
{The $m$-out-of-$n$ bootstrap for heavy tails (paper's Example 4); the
cross-sectional pairs bootstrap for dynamic panels (Example 5). Both are
discussed in the supplement.}

% ----------------
\slide{22}{Corollary 3.1.}
\say{$H(\pnh)\dto U[0,1]$. Proof: by Theorem 3.1, $\pnh\dto V$ where
$V=G(F^{-1}(U))$. Apply $H=F\circ G^{-1}$ to $V$: the inner $G$s cancel,
leaving $F\circ F^{-1}(U)=U$. By CMT, $H(\pnh)\dto U$.

The reading is that applying $H$ \emph{unwarps} the distorted $\pnh$ back to
uniform --- exactly what we wanted. Of course, $H$ is unknown, so we need
$\Hnh$.}

\trans{Two ways to estimate $H$: plug-in, then double bootstrap.}

% ----------------
\slide{23}{Approach A --- Plug-in (Corollary 3.2).}
\say{Assume $H=H_\gamma$ for some finite-dimensional $\gamma$, with $\hat\gamma_n
\pto\gamma$. Corollary 3.2: if $H_\gamma(u)$ is jointly continuous in $(\gamma,u)$,
then $\pnt:=H_{\hat\gamma_n}(\pnh)\dto U[0,1]$.

The proof is two pieces. First, by joint continuity in $\gamma$,
$\sup_u|H_{\hat\gamma_n}(u)-H_\gamma(u)|\pto 0$, so $\pnt-H(\pnh)\pto 0$.
Second, $H(\pnh)\dto U[0,1]$ by Corollary 3.1. Slutsky combines them.

The plug-in approach is cheap --- one bootstrap plus one formula --- but
requires deriving $H_\gamma$ in closed form. In the Gaussian case, which
covers our examples, $H$ takes a simple shape; that is the next slide.}

\trans{Gaussian case --- the variance-ratio correction.}

\textbf{Anticipated questions:}
\Q{If $\hat\gamma_n$ is only $n^{1/2}$-consistent, do we need any rate of convergence?}
{No, only consistency in probability. The proof goes through Slutsky;
no rate is required.}

% ----------------
\slide{24}{Gaussian plug-in --- the correction $m$.}
\say{Suppose $\xi_1\sim N(0,v^2)$ and $(\xi_1,\xi_2)$ are jointly normal. Write
$\xi_1-\xi_2\sim N(0,v_d^2)$. Then $G(u)=\Phi(u/v)$, $F(u)=\Phi(u/v_d)$, and
\[
  H(u) \;=\; \Phi\bigl(m^{-1}\Phi^{-1}(u)\bigr),
  \qquad m \;=\; v_d/v.
\]
\emph{Reading.} Without (1.3), $\Tn-\Bnh\dto N(0,v^2)$ --- the bootstrap spread
is $v^2$. Under (1.3), the actual spread is $v_d^2 \neq v^2$. The correction
$\pnt=\Phi(m^{-1}\Phi^{-1}(\pnh))$ rescales the $p$-value as if the spread had
always been $v_d^2$. \emph{$B$ drops out completely} --- only standard deviations
appear.

In our two examples: ridge has $m^2$ depending on $\Sigma_{xx}$ and $c_0$, but
\emph{not} on the unknown local-to-zero $\delta$ --- so $m$ is estimable.
Nonparametric is even better: $\sigma^2$ cancels and $m$ depends \emph{only on
the kernel $K$} --- a known constant, no estimation at all.}

\textbf{Anticipated questions:}
\Q{Why does $\sigma^2$ cancel only in nonparametric?}
{In the nonparametric case both $v$ and $v_d$ scale linearly with $\sigma^2$,
so the ratio is dimensionless in $\sigma$. In ridge the correlation between
$\xi_1$ and $\xi_2$ is different and $\sigma^2$ does not cancel out.}
\Q{Is $m$ always less than 1, or always greater?}
{Generally $m\ge 1$ in these examples: $\xi_1-\xi_2$ has more spread than
$\xi_1$ because of the extra randomness from $\xi_2$. So the correction
shrinks $\Phi^{-1}(\pnh)$ toward zero, making the test more conservative.}
\Q{Why is this a ``ratio of standard deviations'' rather than something with $B$?}
{Because both $G$ and $F$ are cdfs of \emph{centred} variables; $B$ has been
subtracted away by definition. So $H$ depends only on the variance structure
of $(\xi_1,\xi_2)$, not on the bias.}

% ----------------
\slide{25}{Approach B --- Double bootstrap.}
\say{If you cannot find $H_\gamma$ analytically, use a double bootstrap. From
each first-level sample $\Dns$, draw a second-level $\Dnss$. Compute the
second-level $p$-value $\pnhs$. Estimate $H$ by the bootstrap cdf of $\pnhs$.
$\pnt$ is then $\Pstar(\pnhs\le\pnh)$ --- the fraction of second-level
$p$-values below the original.

Theorem 3.2 says this works under Assumptions 1--3. Assumption 3 just says the
double bootstrap replicates both the noise (at level 2) and the bias gap.}

\trans{Corollary 3.3 gives a much cheaper alternative.}

\textbf{Anticipated questions:}
\Q{Computational cost?}
{$B^2$ where $B$ is the bootstrap size --- typically prohibitive. That is why
Corollary 3.3 matters.}

% ----------------
\slide{26}{Corollary 3.3 --- the practical shortcut.}
\say{The double bootstrap modified $p$-value is \emph{asymptotically equivalent}
to a single-level bootstrap applied to the bias-corrected statistic
$\Tn-\Bnh$. So if $\Bnh$ is available in closed form (model averaging, ridge),
no double bootstrap is needed: just bootstrap the bias-corrected statistic.

Note: $\Bnh$ does \emph{not} need to be a consistent estimator of $\Bn$.
We are not estimating bias --- we are using whatever the bootstrap provides
as the centring device.}

\trans{Examples revisited.}

\textbf{Anticipated questions:}
\Q{Doesnt this look like just bias-correcting?}
{No. We are not bias-correcting in the sense of plugging in an estimator
of $B$. We are using $\Bnh$ as a centring; $\Bnh$ is biased for $\Bn$ in
general (that is the whole point). The validity comes from Theorem 3.2,
not from $\Bnh\to B$.}

% ====================================================================
\section{Part 4 --- Examples revisited (slides 27--28)}

\slide{27}{Ridge --- implementing the fix.}
\say{Same recipe. The variance ratio simplifies to a ratio involving
$\Sigma_{xx}$ and the regularised version $\tilde\Sigma_{xx}$.
Sanity check: when $c_0=0$ (no shrinkage), $m=1$ and there is no correction ---
ridge collapses to OLS, which is unbiased.}

% ----------------
\slide{28}{Nonparametric --- implementing the fix.}
\say{The cleanest case: $m$ depends \emph{only on the kernel}, no estimation
needed. Compute $\hat\beta_h(x)$ at the MSE-optimal bandwidth, run the
parametric bootstrap, look up $m$ from the kernel, report
$\pnt=\Phi(m^{-1}\Phi^{-1}(\pnh))$.

Compare to existing approaches. Undersmoothing kills the bias but at the cost
of power. Robust bias correction (CCF 2018) needs a second bandwidth. This
paper: keep the MSE-optimal $h$, one-line correction, no estimation of
$\beta''(x)$. Cleaner and finite-sample friendly.}

\trans{Wrap-up.}

\textbf{Anticipated questions:}
\Q{Why does the paper not include simulations?}
{They are in the supplement.}

% ====================================================================
\section{Part 5 --- Wrap-up}

\slide{29}{Summary.}
\say{Four bullets.

\emph{Problem.} Asymptotic bias $B$ that the bootstrap cannot replicate
makes $\pnh$ non-uniform.

\emph{Diagnosis.} The bootstrap bias $\Bnh$ misses $B$ by a non-vanishing random
$\xi_2$. Source: an inestimable quantity at the $n^{-1/2}$ boundary.

\emph{Fix.} The asymptotic cdf $H=F\circ G^{-1}$ of $\pnh$ is continuous and
free of $B$. So $\pnt=\Hnh(\pnh)\dto U[0,1]$.

\emph{Two routes.} Plug-in (Gaussian: variance ratio $m=v_d/v$) or double
bootstrap. Corollary 3.3 reduces the latter to a single bootstrap of
$\Tn-\Bnh$.

The bigger picture for the reading group: a third role for the bootstrap.
Not variance, not refinement, but \emph{rescuing inference} when the
first-order bootstrap is invalid --- without ever estimating the bias.}

\trans{Questions.}

% ====================================================================
\section{Catch-all questions to be ready for}

\Q{Why not just do uniform inference using subsampling (Politis--Romano)?}
{Subsampling is consistent under weaker conditions, but it can be very
imprecise in finite samples and choosing the subsample size $b$ is delicate.
Prepivoting requires stronger assumptions but, when applicable, gives a
clean closed-form correction (especially in the Gaussian case).}

\Q{Is this the same as Andrews--Guggenberger size-correction?}
{Both target asymptotic non-uniformity, but the mechanisms differ.
Andrews--Guggenberger correct for non-uniform convergence over the
parameter space; Cavaliere et al.\ correct for a random bootstrap distribution
in the limit at a single parameter value.}

\Q{What happens under heteroskedasticity / dependence?}
{Use a bootstrap that handles it (wild, block) and re-derive Assumption 2.
The framework does not change; only the verification of (1.2) and (1.3).}

\Q{Could $H$ depend on a nuisance parameter that is itself not consistently estimable?}
{Yes, in principle. The plug-in approach then fails, but the double bootstrap
can still work --- it estimates $H$ nonparametrically. This is exactly the
selling point of double bootstrap.}

\Q{Why prepivoting and not simply bootstrap-of-bootstrap on the studentised statistic?}
{Studentisation only fixes scale, not bias. The non-uniformity here is
location-driven (random $\Bnh-B$). Studentisation does not address that.}

\Q{If $\Bnh$ is implicit, how do we even compute the correction in practice?}
{For Gaussian limits we never compute $\Bnh$ directly; we estimate the
\emph{variances} $v^2$ and $v_d^2$. For Corollary 3.3 we do need a closed
form for $\Bnh$, which the paper provides for ridge and model averaging.}

\Q{How robust is the prepivoted test to misspecification?}
{Same robustness as the underlying bootstrap. Prepivoting is a transformation;
it does not strengthen distributional assumptions but inherits them.}




% ============================================================
\section{Modes of Convergence Used in the Paper}
% ============================================================
 
\subsection{Standard convergence}
 
\begin{definition}[Convergence in probability]
  $Y_n \xrightarrow{p} \xi$ means: for every $\varepsilon > 0$,
  \[
    P(|Y_n - \xi| > \varepsilon) \to 0 \quad\text{as } n \to \infty.
  \]
  The random variable $Y_n$ gets arbitrarily close to $\xi$ with
  probability approaching~1. When $\xi = c$ is a constant, this is the
  usual consistency statement (e.g.\ $\hat\sigma^2_n \xrightarrow{p} \sigma^2$).
\end{definition}
 
\begin{definition}[Convergence in distribution]
  $Y_n \xrightarrow{d} \xi$ means: for every continuity point $u$ of
  the cdf of $\xi$,
  \[
    P(Y_n \le u) \to P(\xi \le u).
  \]
  This is weaker than convergence in probability: $Y_n$ and $\xi$ need not
  be close in any single realisation; only their \textit{distributions}
  converge. The CLT is the canonical example:
  $\sqrt{n}(\bar X_n - \mu)/\sigma \xrightarrow{d} N(0,1)$.
\end{definition}
 
% ============================================================
\subsection{Bootstrap convergence}
 
\begin{definition}[Bootstrap weak convergence in probability: $\xrightarrow{d^*}_p$]
  $Y_n^* \xrightarrow{d^*}_p \xi$ means: for every continuity point $u$ of
  the cdf $G(u) = P(\xi \le u)$,
  \[
    P^*(Y_n^* \le u) - G(u) \xrightarrow{p} 0.
  \]
  Here $P^*(\cdot) = P(\cdot \mid D_n)$ is the \textbf{conditional}
  probability given the original data~$D_n$. The quantity
  $P^*(Y_n^* \le u)$ is itself a random variable (it depends on which
  $D_n$ was observed). The statement says: this random cdf converges
  \textit{in probability} (over realisations of $D_n$) to the
  deterministic cdf $G$.
 
  \medskip
 
  \textbf{Why ``in probability''?}
  Different realisations of $D_n$ give different bootstrap distributions.
  For most datasets the bootstrap cdf will be close to $G$, but for some
  pathological datasets it might not be. The ``in probability'' qualifier
  says the set of bad datasets has vanishing probability.
 
  \medskip
 
  \textbf{Equivalently:} $Y_n^* = o_{p^*}(1)$ means
  $P^*(|Y_n^*| > \varepsilon) \xrightarrow{p} 0$ for all $\varepsilon > 0$.
\end{definition}
 
\begin{definition}[Double bootstrap convergence: $\xrightarrow{d^{**}}_{p^*}$, in probability]
  $Y_n^{**} \xrightarrow{d^{**}}_{p^*} \xi$, in probability, means: for
  every continuity point $u$ of $G$,
  \[
    P^{**}(Y_n^{**} \le u) - G(u) \xrightarrow{p^*}_p 0.
  \]
  Now $P^{**}(\cdot) = P(\cdot \mid D_n^*, D_n)$ is conditional on both
  the first-level bootstrap data $D_n^*$ and the original data $D_n$.
  The convergence $\xrightarrow{p^*}_p 0$ means: the $P^*$-probability
  that $|P^{**}(Y_n^{**} \le u) - G(u)|$ exceeds~$\varepsilon$ converges
  to zero in $P$-probability.
 
  \medskip
 
  In plain language: for most original datasets $D_n$, and for most
  first-level bootstrap datasets $D_n^*$ drawn from $D_n$, the
  second-level bootstrap cdf is close to $G$.
\end{definition}
 
% ============================================================
\subsection{How they relate in the paper}
 
\begin{center}
\begin{tabular}{@{}lll@{}}
  \toprule
  Statement & What's random & Conditioned on \\
  \midrule
  $T_n - B_n \xrightarrow{d} \xi_1$
    & $T_n$ (through $D_n$)
    & nothing \\[4pt]
  $T_n^* - \hat B_n \xrightarrow{d^*}_p \xi_1$
    & $T_n^*$ (through bootstrap variates)
    & $D_n$ \\[4pt]
  $T_n^{**} - \hat B_n^* \xrightarrow{d^{**}}_{p^*} \xi_1$
    & $T_n^{**}$ (through level-2 variates)
    & $D_n^*$ and $D_n$ \\
  \bottomrule
\end{tabular}
\end{center}
 
Each level adds one layer of conditioning. The ``in probability''
qualifiers handle the fact that the conditioning set is itself random.
 
% ============================================================
\section{Two Key Theorems Used in the Proofs}
% ============================================================
 
\subsection{Polya's theorem}
 
\begin{remark}[Polya's theorem]
  If $F_n(u) \to F(u)$ pointwise and $F$ is \textbf{continuous}, then
  the convergence is \textbf{uniform}:
  \[
    \sup_{u \in \R} |F_n(u) - F(u)| \to 0.
  \]
\end{remark}
 
\textbf{Why the paper needs this.}
In the proof of Theorem~3.1, we have
$P^*(T_n^* - \hat B_n \le u) \to G(u)$ pointwise (from Assumption~2(i)).
We then evaluate this at the \textit{random} point $u = T_n - \hat B_n$.
Pointwise convergence alone does not let you plug in a random argument ---
the error could be large exactly where the random point lands.
Uniform convergence (via Polya, using continuity of $G$) guarantees the
error is small \textit{everywhere simultaneously}, so plugging in a
random point is safe:
\[
  |\hat p_n - G(T_n - \hat B_n)|
  \le \sup_u |P^*(T_n^* - \hat B_n \le u) - G(u)|
  \xrightarrow{p} 0.
\]
 
\textbf{This is why Assumption~1 requires $G$ continuous.}
Without continuity, Polya does not apply, uniform convergence fails,
and the proof breaks at Step~1.
 
\subsection{The probability integral transform (PIT)}
 
\begin{remark}[PIT]
  If $X$ has continuous cdf $F$, then $F(X) \sim U[0,1]$.
\end{remark}
 
\textit{Proof in one line:}
$P(F(X) \le u) = P(X \le F^{-1}(u)) = F(F^{-1}(u)) = u$.
 
\medskip
 
The paper uses the PIT twice:
\begin{enumerate}[label=(\arabic*)]
  \item In the \textbf{benchmark} (consistency case):
    $G(\xi_1) \sim U[0,1]$ because $\xi_1$ has cdf $G$.
  \item In \textbf{Corollary~3.1}:
    $H(\hat p_n) \xrightarrow{d} U[0,1]$ because $\hat p_n$ has
    asymptotic cdf $H$.
\end{enumerate}
 
\textbf{This is why $G$ must be strictly increasing} (for use~1: ensures
$G^{-1}$ exists cleanly) and \textbf{$H$ must be continuous} (for use~2:
ensures the transform works for the $p$-value).
 
% ============================================================
\section{Ridge Regression Example --- Full Walkthrough}
% ============================================================
 
\subsection{Why ridge? Where is it used?}
 
Ridge regression (and penalised estimation more broadly: LASSO, elastic
net) is the standard tool when you have many regressors with limited
individual explanatory power. You shrink coefficients toward zero to
reduce variance at the cost of some bias. This is now ubiquitous in
applied econometrics --- high-dimensional controls in causal inference,
machine learning for prediction, regularised IV, etc.
 
\subsection{The model}
 
Linear regression $y_t = \theta'x_t + \varepsilon_t$, $t = 1,\ldots,n$,
where $x_t$ is a $p \times 1$ non-stochastic vector and
$\varepsilon_t \sim \text{iid}(0,\sigma^2)$.
 
\medskip
 
Interest is in testing $H_0 : g'\theta = r$ for a known vector $g$
(e.g.\ $g = e_j$ tests the $j$-th coefficient).
 
\subsection{The ridge estimator}
 
The ridge estimator minimises the penalised least squares criterion:
\[
  \frac{1}{n}\|y - X\theta\|^2 + \frac{c_n}{n}\|\theta\|^2
\]
giving the closed-form solution:
\[
  \tilde\theta_n
  = \Bigl(S_{xx} + \frac{c_n}{n}I_p\Bigr)^{-1} S_{xy}
  =: \tilde S_{xx}^{-1} S_{xy}
\]
where $S_{xx} = X'X/n$ and $S_{xy} = X'y/n$.
 
\medskip
 
\textbf{What does ridge do mechanically?}
OLS inverts $S_{xx}$; ridge inverts $S_{xx} + \frac{c_n}{n}I_p$, a
\textit{larger} matrix (in the positive-definite sense). Inverting a
larger matrix gives smaller coefficients --- ridge systematically pulls
$\tilde\theta_n$ toward zero. When $\theta \ne 0$, this pull is bias.
 
\subsection{The source of bias}
 
The key algebraic identity:
\[
  \tilde S_{xx}^{-1} S_{xx} - I_p = -\frac{c_n}{n}\tilde S_{xx}^{-1}
\]
This follows from $S_{xx} = \tilde S_{xx} - \frac{c_n}{n}I$, so
$\tilde S_{xx}^{-1} S_{xx} = I - \frac{c_n}{n}\tilde S_{xx}^{-1}$.
 
\medskip
 
Using $S_{xy} = S_{xx}\theta + X'\varepsilon/n$:
\[
  \tilde\theta_n - \theta
  = \underbrace{-\frac{c_n}{n}\tilde S_{xx}^{-1}\theta}_{\text{bias (shrinkage)}}
  + \underbrace{\tilde S_{xx}^{-1}\frac{X'\varepsilon}{n}}_{\text{noise}}
\]
 
The bias is proportional to $\frac{c_n}{n}$ (penalty strength) times
$\tilde S_{xx}^{-1}\theta$ (shrinkage direction). This is entirely
mechanical --- ridge pulls $\tilde\theta_n$ toward zero, and the amount
of pull depends on how large $\theta$ is and how strong the penalty is.
 
\subsection{The local-to-zero framework}
 
Set $\theta = \delta n^{-1/2}$ and $c_n/n \to c_0 > 0$.
 
\medskip
 
\textbf{Why this parametrisation?}
\begin{itemize}
  \item If $\theta$ is fixed: the normalised bias
    $n^{1/2} \cdot \frac{c_n}{n}\tilde S_{xx}^{-1}\theta \to \infty$.
    Ridge is terrible. Not interesting.
  \item If $\theta = 0$: no bias. Ridge $=$ OLS asymptotically.
    Not interesting either.
  \item If $\theta = \delta n^{-1/2}$: the normalised bias
    $\frac{c_n}{n}\tilde S_{xx}^{-1}\delta \to c_0\tilde\Sigma_{xx}^{-1}\delta$,
    a nonzero constant. \textit{This is the boundary where bias and noise
    are the same asymptotic order} --- the regime where shrinkage
    genuinely helps but also genuinely distorts inference.
\end{itemize}
 
\subsection{The test statistic and (1.1)}
 
Multiply by $n^{1/2}$ and apply $g'$:
\[
  T_n := n^{1/2}(g'\tilde\theta_n - r)
  = \underbrace{-\frac{c_n}{n}g'\tilde S_{xx}^{-1}\delta}_{B_n \;\to\; B \,=\, -c_0 g'\tilde\Sigma_{xx}^{-1}\delta}
  + \underbrace{g'\tilde S_{xx}^{-1}\frac{X'\varepsilon}{n^{1/2}}}_{\xi_{1,n} \;\xrightarrow{d}\; N(0, v^2)}
\]
where $v^2 = \sigma^2 g'\tilde\Sigma_{xx}^{-1}\Sigma_{xx}\tilde\Sigma_{xx}^{-1}g$.
 
\medskip
 
This confirms (1.1): $T_n - B_n \xrightarrow{d} \xi_1 \sim N(0,v^2)$
with $B = -c_0 g'\tilde\Sigma_{xx}^{-1}\delta \ne 0$.
 
\subsection{Why $B$ cannot be consistently estimated}
 
To estimate $B = -c_0 g'\tilde\Sigma_{xx}^{-1}\delta$, we can estimate
$c_0$, $\tilde\Sigma_{xx}$, and $g$ consistently. But $\delta$ requires:
\[
  n^{1/2}\hat\theta_n = n^{1/2}\theta + n^{1/2}(\hat\theta_n - \theta)
  = \delta + O_p(1).
\]
The $O_p(1)$ noise never vanishes --- we observe $\delta$ plus
irreducible noise, and no amount of data can separate them.
 
\subsection{The bootstrap and (1.2)}
 
Pairs bootstrap: resample $(y_t^*, x_t^*)$ iid from $(y_t, x_t)$.
Compute $\tilde\theta_n^* = \tilde S_{x^*x^*}^{-1}S_{x^*y^*}$ and
\[
  T_n^* = n^{1/2}g'(\tilde\theta_n^* - \hat\theta_n).
\]
 
\textbf{Why centre at $\hat\theta_n$ (OLS), not $\tilde\theta_n$ (ridge)?}
In the bootstrap world, the ``true parameter'' is whatever makes
bootstrap errors uncorrelated with regressors. OLS satisfies
$\sum x_t(y_t - x_t'\hat\theta_n) = 0$, so
$E^*[x_t^*\varepsilon_t^*] = 0$. Ridge does not satisfy this first-order
condition --- centring at $\tilde\theta_n$ would create spurious
correlation. Same principle as GMM bootstrap (Horowitz).
 
\medskip
 
The implicit bootstrap bias is
$\hat B_n = -\frac{c_n}{n}g'\tilde S_{xx}^{-1}n^{1/2}\hat\theta_n$
(using $\tilde S_{x^*x^*}^{-1} - \tilde S_{xx}^{-1} = o_{p^*}(1)$).
Then:
\[
  T_n^* - \hat B_n \xrightarrow{d^*}_p N(0, v^2).
\]
This confirms (1.2): the bootstrap captures the noise correctly.
 
\subsection{The failure: (1.3)}
 
\[
  \hat B_n - B_n
  = -\frac{c_n}{n}g'\tilde S_{xx}^{-1} \cdot n^{1/2}(\hat\theta_n - \theta)
  \xrightarrow{d} \xi_2 \sim N(0, v_{22}), \quad v_{22} > 0.
\]
This holds because $n^{1/2}(\hat\theta_n - \theta) = O_p(1)$, not
$o_p(1)$. The bootstrap bias $\hat B_n$ tracks $B_n$ with irreducible
noise of the same order as $B_n$ itself.
 
\medskip
 
\textbf{Key structural insight:} both $T_n - B_n$ and $\hat B_n - B_n$
are linear in the same score $n^{1/2}S_{x\varepsilon}$:
\[
  \begin{pmatrix}T_n - B_n \\ \hat B_n - B_n\end{pmatrix}
  = \begin{pmatrix}g'\tilde S_{xx}^{-1} \\
    -c_n n^{-1}g'\tilde S_{xx}^{-1}S_{xx}^{-1}\end{pmatrix}
  n^{1/2}S_{x\varepsilon}
  \xrightarrow{d} N(0, V)
\]
They are jointly normal and correlated because they share the same
source of randomness.
 
\subsection{The fix}
 
Lemma~4.4 verifies Assumptions~1 and 2 with $(\xi_1, \xi_2) \sim N(0,V)$.
The distortion parameter simplifies to:
\[
  m^2 = \frac{g'\Sigma_{xx}^{-1}g}
             {g'\tilde\Sigma_{xx}^{-1}\Sigma_{xx}\tilde\Sigma_{xx}^{-1}g}
\]
using the algebraic identity
$\Sigma_{xx} + 2c_0 I + c_0^2\Sigma_{xx}^{-1}
= \tilde\Sigma_{xx}\Sigma_{xx}^{-1}\tilde\Sigma_{xx}$.
 
\medskip
 
\textbf{Plug-in recipe:}
\begin{enumerate}
  \item Compute ridge $\tilde\theta_n$ and
    $T_n = n^{1/2}(g'\tilde\theta_n - r)$.
  \item Run pairs bootstrap: resample $(y_t^*, x_t^*)$, centre at OLS:
    $T_n^{*(b)} = n^{1/2}g'(\tilde\theta_n^{*(b)} - \hat\theta_n)$.
    Compute $\hat p_n = B^{-1}\sum_b \mathbf{1}\{T_n^{*(b)} \le T_n\}$.
  \item Estimate $m$ from sample moments:
    $\hat m_n^2 = (g'S_{xx}^{-1}g) / (g'\tilde S_{xx}^{-1}S_{xx}\tilde S_{xx}^{-1}g)$.
  \item Report $\tilde p_n = \Phi(\hat m_n^{-1}\Phi^{-1}(\hat p_n))$.
\end{enumerate}
 
\medskip
 
\textbf{Sanity checks:}
\begin{itemize}
  \item $c_0 = 0 \Rightarrow \tilde\Sigma_{xx} = \Sigma_{xx}
    \Rightarrow m = 1$: no correction needed (ridge $=$ OLS).
  \item $m$ depends on $\Sigma_{xx}$ and $c_0$ (estimable), \textit{not}
    on $\delta$ (inestimable). The unknown bias $B$ drops out completely.
  \item $\hat m_n \xrightarrow{p} m$ because $S_{xx} \to \Sigma_{xx}$ and
    $\tilde S_{xx} \to \tilde\Sigma_{xx}$.
\end{itemize}
 
Lemma~4.5 verifies Assumption~3, so the double bootstrap is also valid.
 

\end{document}
